\documentclass[10pt,a4paper]{amsart}
\usepackage[margin=19mm]{geometry}
\usepackage{amsmath,amssymb,mathtools}
\usepackage[scr=rsfs,cal=euler]{mathalfa}
\usepackage{xfrac}
\usepackage[unicode,hidelinks,bookmarks=false]{hyperref}
\newcommand{\ie}{i.e.\ }
\newcommand{\cf}{cf.\ }
\newcommand{\resp}{respectively\ }
\newcommand{\Subsection}{\subsection}
\providecommand{\nobreakdash}{\nobreak\hbox{-}}
\setlength{\emergencystretch}{2em}
\multlinegap0pt
\usepackage{mathrsfs}
\usepackage{appendix}
\usepackage{fancyhdr}
\usepackage{charter}
\usepackage{typearea}
\usepackage{color}
\usepackage{mathrsfs}
\renewcommand{\baselinestretch}{1.10}
\usepackage{enumitem}
\usepackage{color}
\newcommand{\comment}[1]{}
      \def\@setcopyright{}
      \def\serieslogo@{}
\newcommand{\Complex}{\mathbb C}
\newcommand{\Real}{\mathbb R}
\newcommand{\N}{\mathbb N}
\newcommand{\ddbar}{\overline\partial}
\newcommand{\p}{\partial}
\newcommand{\CC}{\mathbb C}
\newcommand{\pr}{\partial}
\newcommand{\ol}{\overline}
\newcommand{\ov}{\overline}
\newcommand{\sm}{\setminus}
\newcommand{\Td}{\widetilde}
\newcommand{\norm}[1]{\left\Vert#1\right\Vert}
\newcommand{\abs}[1]{\left\vert#1\right\vert}
\newcommand{\rabs}[1]{\left.#1\right\vert}
\newcommand{\set}[1]{\left\{#1\right\}}
\newcommand{\To}{\rightarrow}
\DeclareMathOperator{\Ker}{Ker}
\newcommand{\pit}{\mathit\Pi}
\newcommand{\cali}[1]{\mathscr{#1}}
\newcommand{\cH}{\cali{H}}
\newcommand{\mU}{\mathcal{U}}
\newcommand{\mV}{\mathcal{V}}
\def\cC{\mathscr{C}}
\def\cL{\mathscr{L}}
\newtheorem{theorem}{Theorem}[section]
\newtheorem{lemma}[theorem]{Lemma}
\newtheorem{claim}[theorem]{Claim}
\newtheorem{corollary}[theorem]{Corollary}
\newtheorem{question}[theorem]{Question}
\newtheorem{proposition}[theorem]{Proposition}
\newtheorem{problem}[theorem]{Problem}
\newtheorem{definition}[theorem]{Definition}
\newtheorem*{conjecture}{Conjecture}
\newtheorem{example}[theorem]{Example}
\newtheorem{remark}[theorem]{Remark}
\newtheorem{assumption}[theorem]{Assumption}
\numberwithin{equation}{section}
\setlength{\headheight}{14pt}
\begin{document}
\title[A note on csc Bergman metric]{A note on csc Bergman metric}
\author{Xiaojun Huang; Xiaoshan Li}
\date{}
\maketitle
\noindent\emph{Source and licence.} A note on csc Bergman metric. Version arXiv:2606.01319v1 (CC BY 4.0). This material is distributed under the Creative Commons Attribution 4.0 International licence, http://creativecommons.org/licenses/by/4.0/, as recorded in the arXiv OAI licence metadata for this exact version and shown on the abstract page https://arxiv.org/abs/2606.01319v1. Full rights notice: licenses/arxiv-2606.01319-CC-BY-4.0.txt.\par\smallskip
\noindent\textit{Editorial scope.} Counted unit: the original \texttt{remark} environment \texttt{rem1} at source lines 273--275 together with its complete proof at lines 276--339; the delivered document keeps source lines 270--340 so that every referenced label and every citation resolves inside it. Native editable source is the paper's own concatenated TeX; the AMS wrapper is editorial formatting only. No publisher class, upstream script or unrelated artwork is redistributed.
\section{Remarks}

%We make the following more remarks in this section.

\begin{remark}\label{rem1}
    For $m>1$ an integer  and $n\geq 2$, let $$\mathcal E_m:=\left\{z\in\mathbb C^{n-1}\times \mathbb C:\sum_{j=1}^{n-1}|z_j|^2+|z_n|^{2m}<1\right\}.$$ Then we claim that the Bergman metric of $\mathcal E_m$ can not have constant scaler curvature.
\end{remark}

\begin{proof}[Proof of the statement in Remark \ref{rem1}]
    We denote by  $S$  the scalar curvature of the Bergman metric of $\mathcal E_m$. If $\mathcal E_m$ has a csc Bergman metric, then $S\equiv -n$ on $\mathcal E_m$ by \cite[Corollary 2]{KYu96} as it has limit $-n$ at any strongly pseudoconvex point. From \cite {KYu96}, 
    \begin{equation}\label{scaler}
    S(0)=n(n+1)-4a_0\sum_{j=1}^n\frac{b_{jj}}{a_j^2}-a_0\sum_{j\neq k}\frac{b_{jk}}{a_ja_k}.
    \end{equation}
    Here, $a_0=\frac{1}{vol(\mathcal E_m)}$, $a_j=\frac{1}{\|z_j\|_{\mathcal E_m}^2}$, $b_{jk}=\frac{1}{\|z_jz_k\|^2_{\mathcal E_m}}$ with $\|\cdot\|$ the $L^2$-norm on $\mathcal E_m$.
    By direct calculations, 
    \begin{equation}\label{5-28-a1}
\begin{aligned}
a_0=\frac{m\Gamma\!\left(n + \frac{1}{m}\right)}{\pi^{n} \; \Gamma\!\left(\frac{1}{m}\right)},~
a_1=\dots=a_{n-1}=\frac{m\Gamma\!\left(n +1+ \frac{1}{m}\right)}{\pi^{n} \; \Gamma\!\left(\frac{1}{m}\right)},~ 
a_n=\frac{m\;\Gamma\!\left(n+\frac{2}{m}\right)}{\pi^{n} \; \Gamma\!\left(\frac{2}{m}\right)}
\end{aligned}
\end{equation}
and 
\begin{equation}\label{5-28-a2}
\begin{split}
&b_{jj}=\dfrac{m}{2\pi^n}\,\dfrac{\Gamma\!\left(n+2+\frac1m\right)}{\Gamma\!\left(\frac1m\right)},~ j=1,\dots,n-1; ~b_{nn}=\dfrac{m}{\pi^n}\,\dfrac{\Gamma\!\left(n+\frac3m\right)}{\Gamma\!\left(\frac3m\right)},\\
&b_{jk}=\dfrac{m}{\pi^n}\,\dfrac{\Gamma\!\left(n+2+\frac1m\right)}{\Gamma\!\left(\frac1m\right)},~  j\neq k,\; j,k=1,\dots,n-1,\\
&b_{jn}=b_{nj}=\dfrac{m}{\pi^n}\,\dfrac{\Gamma\!\left(n+1+\frac2m\right)}{\Gamma\!\left(\frac2m\right)},  j=1,\dots,n-1
\end{split}
 \end{equation}   
Substituting (\ref{5-28-a1}) and (\ref{5-28-a2}) to (\ref{scaler}) we have
\begin{equation}
S(0)=2-\frac{(n-1)(n+\frac 2m)}{n+\frac 1m}
-4\,\frac{\Gamma\!\left(n+\frac1m\right)\,\
\Gamma\!\left(n+\frac3m\right)\Gamma\!\left(\frac2m\right)^2\,}{\Gamma\!\left(\frac1m\right)\,\Gamma\!\left(\frac3m\right)\,
\Gamma\!\left(n+\frac2m\right)^2}.
\end{equation}
Set $a=\frac{1}{m}\in(0, 1)$, $H_n(a)
=
\frac{
\Gamma(n+3a)\Gamma(n+a)\Gamma(2a)^2
}{
\Gamma(3a)\Gamma(a)\Gamma(n+2a)^2
}=\prod_{j=0}^{n-1}
\frac{(j+a)(j+3a)}{(j+2a)^2}$ and $T_n(a)
=\frac{3n+(4-n)a}{4(n+a)}.$ Then 
\begin{equation}
\begin{split}
    S(0)+n=4(T_n(a)-H_n(a)).
    \end{split}
\end{equation}
By direct calculation $T_2(a)<H_2(a)$.  Since
$H_{n+1}(a)=H_n(a)
\frac{(n+a)(n+3a)}{(n+2a)^2}$ and 
$$
T_n(a)\frac{(n+a)(n+3a)}{(n+2a)^2}-T_{n+1}(a)=\frac{an(1-a)^2}{4(n+2a)^2(n+1+a)}.
$$
Since $0<a<1$ and $n\geq 2$, the right-hand side is strictly positive,
hence
\[T_n(a)\frac{(n+a)(n+3a)}{(n+2a)^2}>T_{n+1}(a).
\]
Assume inductively that
\[H_n(a)>T_n(a).\]
Then
\[H_{n+1}(a)=H_n(a)\frac{(n+a)(n+3a)}{(n+2a)^2}>T_n(a)\frac{(n+a)(n+3a)}{(n+2a)^2}>T_{n+1}(a).\]
Thus
\[H_n(a)>T_n(a)\]
for every $n\ge2$.
Consequently,
\[S(0)+n=4(T_n(a)-H_n(a))<0.\]
Therefore, $S(0)<-n$ and thus we get a contradiction.
\end{proof}


\begin{thebibliography}{99}
\bibitem[KYu96]{KYu96} S. Krantz and J. Yu: On the Bergman invariant and curvatures of the Bergman metric, Illinois J. Math. 40 (1996), no. 2, 226--244. \bibitem [HX21]{HX21} X. Huang and M. Xiao, Bergman-Einstein metrics, a generalization of Kerner's theorem and Stein spaces with spherical boundaries. J. Reine Angew. Math. 770 (2021), 183–203.
\end{thebibliography}
\end{document}
